% GENERATED FILE. Edit canonical lessons, then run npm run generate:books.
% canonical-source-hash: fnv1a64:cbe5d8e42282feb1
% canonical-lessons: SA-C244-ucitah, SA-C244-anucitah, SA-C244-avasyakah, SA-C244-anyah, SA-C244-ekaki

\chapter{Proper, Improper, Necessary, Other, Alone}
\label{ch:sa-proper-improper-necessary-other-alone}
\hlchaptermodality{244}

\begin{chapteropening}
\textbf{By the end of this chapter:} \emph{I can say proper, improper, necessary, other and alone.}

Complete the last lesson of chapter 244: I can say proper, improper, necessary, other and alone.
\end{chapteropening}

\section[\texorpdfstring{\sk{उचितः} (ucitaḥ)}{उचितः (ucitaḥ)}]{\sk{उचितः} (ucitaḥ) — proper, fitting}

\label{lesson:SA-C244-ucitah}



\begin{quote}
Before the new one: say the Sanskrit for happy, then the Sanskrit for unhappy.
\end{quote}

\subsection*{You'll want to know: \sk{उचितः}}
\textbf{\sk{उचितः}} — \emph{ucitaḥ} — \textquotedblleft{}proper, fitting\textquotedblright{}.

\begin{grammarlens}[title={What you've built}]
One more word for this chapter.
\end{grammarlens}

\subsection*{Your turn}
\begin{itemize}
\raggedright
  \item \emph{Say it:} \emph{ucitaḥ}
  \item \emph{Say it:} \emph{ucitaḥ}, once more
  \item \emph{Say it:} say \emph{ucitaḥ}
  \item \emph{Recall:} say the Sanskrit for happy, then the Sanskrit for unhappy, then say \emph{ucitaḥ} again
\end{itemize}

\begin{tcolorbox}[breakable,colback=teal!4,colframe=teal!35!black,title={Before you move on}]
What does \sk{उचितः} mean? (\textquotedblleft{}Proper, fitting\textquotedblright{}.) Say it once more.
\end{tcolorbox}
\section[\texorpdfstring{\sk{अनुचितः} (anucitaḥ)}{अनुचितः (anucitaḥ)}]{\sk{अनुचितः} (anucitaḥ) — improper}

\label{lesson:SA-C244-anucitah}



\begin{quote}
Before the new one: say the Sanskrit for unhappy, then the Sanskrit for proper.
\end{quote}

\subsection*{You'll want to know: \sk{अनुचितः}}
\textbf{\sk{अनुचितः}} — \emph{anucitaḥ} — \textquotedblleft{}improper\textquotedblright{}.

\begin{grammarlens}[title={What you've built}]
One more word for this chapter.
\end{grammarlens}

\subsection*{Your turn}
\begin{itemize}
\raggedright
  \item \emph{Say it:} \emph{anucitaḥ}
  \item \emph{Say it:} \emph{anucitaḥ}, once more
  \item \emph{Say it:} say \emph{anucitaḥ}
  \item \emph{Recall:} say the Sanskrit for unhappy, then the Sanskrit for proper, then say \emph{anucitaḥ} again
\end{itemize}

\begin{tcolorbox}[breakable,colback=teal!4,colframe=teal!35!black,title={Before you move on}]
What does \sk{अनुचितः} mean? (\textquotedblleft{}Improper\textquotedblright{}.) Say it once more.
\end{tcolorbox}
\section[\texorpdfstring{\sk{आवश्यकः} (āvaśyakaḥ)}{आवश्यकः (āvaśyakaḥ)}]{\sk{आवश्यकः} (āvaśyakaḥ) — necessary}

\label{lesson:SA-C244-avasyakah}



\begin{quote}
Before the new one: say the Sanskrit for proper, then the Sanskrit for improper.
\end{quote}

\subsection*{You'll want to know: \sk{आवश्यकः}}
\textbf{\sk{आवश्यकः}} — \emph{āvaśyakaḥ} — \textquotedblleft{}necessary\textquotedblright{}.

\begin{grammarlens}[title={What you've built}]
One more word for this chapter.
\end{grammarlens}

\subsection*{Your turn}
\begin{itemize}
\raggedright
  \item \emph{Say it:} \emph{āvaśyakaḥ}
  \item \emph{Say it:} \emph{āvaśyakaḥ}, once more
  \item \emph{Say it:} say \emph{āvaśyakaḥ}
  \item \emph{Recall:} say the Sanskrit for proper, then the Sanskrit for improper, then say \emph{āvaśyakaḥ} again
\end{itemize}

\begin{tcolorbox}[breakable,colback=teal!4,colframe=teal!35!black,title={Before you move on}]
What does \sk{आवश्यकः} mean? (\textquotedblleft{}Necessary\textquotedblright{}.) Say it once more.
\end{tcolorbox}
\section[\texorpdfstring{\sk{अन्यः} (anyaḥ)}{अन्यः (anyaḥ)}]{\sk{अन्यः} (anyaḥ) — other}

\label{lesson:SA-C244-anyah}



\begin{quote}
Before the new one: say the Sanskrit for improper, then the Sanskrit for necessary.
\end{quote}

\subsection*{You'll want to know: \sk{अन्यः}}
\textbf{\sk{अन्यः}} — \emph{anyaḥ} — \textquotedblleft{}other\textquotedblright{}.

\begin{grammarlens}[title={What you've built}]
One more word for this chapter.
\end{grammarlens}

\subsection*{Your turn}
\begin{itemize}
\raggedright
  \item \emph{Say it:} \emph{anyaḥ}
  \item \emph{Say it:} \emph{anyaḥ}, once more
  \item \emph{Say it:} say \emph{anyaḥ}
  \item \emph{Recall:} say the Sanskrit for improper, then the Sanskrit for necessary, then say \emph{anyaḥ} again
\end{itemize}

\begin{tcolorbox}[breakable,colback=teal!4,colframe=teal!35!black,title={Before you move on}]
What does \sk{अन्यः} mean? (\textquotedblleft{}Other\textquotedblright{}.) Say it once more.
\end{tcolorbox}
\section[\texorpdfstring{\sk{एकाकी} (ekākī)}{एकाकी (ekākī)}]{\sk{एकाकी} (ekākī) — alone}

\label{lesson:SA-C244-ekaki}



\begin{quote}
Before the new one: say the Sanskrit for necessary, then the Sanskrit for other.
\end{quote}

\subsection*{You'll want to know: \sk{एकाकी}}
\textbf{\sk{एकाकी}} — \emph{ekākī} — \textquotedblleft{}alone\textquotedblright{}.

\begin{grammarlens}[title={What you've built}]
One more word for this chapter.
\end{grammarlens}

\subsection*{Your turn}
\begin{itemize}
\raggedright
  \item \emph{Say it:} \emph{ekākī}
  \item \emph{Say it:} \emph{ekākī}, once more
  \item \emph{Say it:} say \emph{ekākī}
  \item \emph{Recall:} say the Sanskrit for necessary, then the Sanskrit for other, then say \emph{ekākī} again
\end{itemize}

\begin{tcolorbox}[breakable,colback=teal!4,colframe=teal!35!black,title={Before you move on}]
What does \sk{एकाकी} mean? (\textquotedblleft{}Alone\textquotedblright{}.) Say it once more.
\end{tcolorbox}
